\chapter{几何演化：输运恒等式与度量变分}

前七章中的几何结构都可以看成固定在一个流形上的对象。若对象随时间变化，必须先区分两种变化：一种只是由时间依赖微分同胚造成的“搬运”，另一种才是度量或其他几何数据本身的真实演化。把二者混在一起会产生大量看似不同的材料导数公式。本章从时间依赖拉回的精确微分公式出发，先统一所有输运方程，再对任意度量变分 $\dot g:=\partial_tg$ 推出逆度量、体积、Levi--Civita 联络与 Hodge 星的变化；Ricci 流 只作为这些一般公式的一个特例。

\section{时间依赖微分同胚与精确输运公式}

设 $\Phi_{t,s}:M\to M$ 是由时间依赖向量场 $V_t$ 生成的两参数流：
\begin{equation*}\partial_t\Phi_{t,s}(p)
=V_t\bigl(\Phi_{t,s}(p)\bigr),
\qquad
\Phi_{s,s}=\id().
\end{equation*}
它满足组合律
\[
\Phi_{t,r}\circ\Phi_{r,s}=\Phi_{t,s}.
\]
若 $T_t$ 是任意时间依赖张量场，则时间导数与拉回满足基本恒等式
\begin{equation}\label{eq:math-dg-transport-identity}
\boxed{
\frac{\dd}{\dd t}\bigl(\Phi_{t,s}^*T_t\bigr)
=\Phi_{t,s}^*
\bigl(\partial_tT_t+\mathcal L_{V_t}T_t\bigr).}
\end{equation}
这一定理是真正的几何版材料导数：$\partial_t$ 比较同一坐标点上的显式时间变化，$\mathcal L_{V_t}$ 补上由流动本身造成的搬运变化。

于是一般线性输运方程
\begin{equation*}\partial_tT_t+\mathcal L_{V_t}T_t=S_t
\end{equation*}
有形式精确解
\begin{equation}\label{eq:math-dg-geometric-transport-solution}
\boxed{
\Phi_{t,s}^*T_t
=T_s+\int_s^t\Phi_{\tau,s}^*S_\tau\,\dd\tau.}
\end{equation}
若 $S_t=0$，则 $\Phi_{t,s}^*T_t=T_s$。对微分形式还由 $\dd\Phi^*=\Phi^*\dd$ 得到：闭性、恰当性与 de Rham 上同调类在纯输运中被保持。

\section{任意度量变分的代数部分}

现在固定底层流形，只令 Riemann 度量 $g(t)$ 光滑依赖时间。记度量的时间导数为
\begin{equation*}\dot g:=\partial_tg.
\end{equation*}
它是一个对称二张量；任何具体几何流都只是给出 $\dot g$ 的特殊表达。

逆度量的变化由恒等式 $g^{ik}g_{kj}=\delta^i{}_j$ 决定：
\[
(\partial_tg^{ik})g_{kj}+g^{ik}\dot g_{kj}=0.
\]
乘以 $g^{j\ell}$ 得
\begin{equation}\label{eq:math-dg-inverse-metric-variation}
\boxed{\partial_tg^{i\ell}=-\dot g^{i\ell}.}
\end{equation}
其中 $\dot g^{i\ell}=g^{ik}g^{\ell j}\dot g_{kj}$ 在当前时间的度量下升指标。

体积形式的变化来自行列式。Jacobi 恒等式
\[
\partial_t\det g
=(\det g)\operatorname{tr}(g^{-1}\dot g)
\]
给出
\begin{equation}\label{eq:math-dg-volume-variation}
\boxed{
\partial_t\operatorname{vol}_g
=\frac12\operatorname{tr}_g(\dot g)\,\operatorname{vol}_g.}
\end{equation}
因此任何保持体积的一阶度量变分都必须满足 $\operatorname{tr}_g \dot g=0$；这比在某组坐标中直接对 $\sqrt{\det g}$ 求导更具几何透明性。

\section{Levi--Civita 联络的变分}

每个 $g(t)$ 都有唯一 Levi--Civita 联络 $\nabla^{(t)}$。固定时间无关的向量场 $X,Y$，定义
\[
\mathcal A(X,Y)
:=\partial_t\bigl(\nabla^{(t)}_XY\bigr).
\]
两个联络之差是张量，因此 $\mathcal A$ 对 $X,Y$ 都是函数线性的。它完全由 $\dot g$ 的一阶协变导数确定。

\eqnrefs{eq:math-dg-connection-variation}说明联络变化只含 $\dot g$ 的一阶导数；因此曲率变化必然含 $\dot g$ 的二阶导数。这一点可以直接写成精确的坐标无关公式。令 $\mathcal A=\partial_t\nabla$，则
\begin{equation}\label{eq:math-dg-curvature-variation-abstract}
\boxed{
\partial_tR(X,Y)Z
=(\nabla_X\mathcal A)(Y,Z)
-(\nabla_Y\mathcal A)(X,Z).}
\end{equation}
固定一组时间无关坐标，记 $\mathcal A^\ell{}_{jk}=\partial_t\Gamma^\ell{}_{jk}$，便有
\[
\partial_tR^\ell{}_{kij}
=\nabla_i\mathcal A^\ell{}_{jk}
-\nabla_j\mathcal A^\ell{}_{ik}.
\]
进一步对曲率算子取迹，得到
\begin{equation}\label{eq:math-dg-ricci-variation}
\partial_t\operatorname{Ric}_{jk}
=\nabla_i\mathcal A^i{}_{jk}
-\nabla_j\mathcal A^i{}_{ik}.
\end{equation}
把\eqnrefs{eq:math-dg-christoffel-variation}代入，并对 $\mathcal R=g^{jk}\operatorname{Ric}_{jk}$ 求导，整理后得到标量曲率的一般变分
\begin{equation}\label{eq:math-dg-scalar-curvature-variation}
\boxed{
\partial_t\mathcal R
=-\langle\dot g,\operatorname{Ric}\rangle_g
+\nabla^i\nabla^j\dot g_{ij}
-\Delta_{\mathrm{LB}}(\operatorname{tr}_g\dot g).}
\end{equation}
这个公式仍未指定任何几何流；Ricci 流、Yamabe 流 或纯 Lie 导数变分都只是对 $\dot g$ 的不同代入。

一般曲率变分公式不仅服务于几何流，也直接给出 Riemann 几何中最重要的变分泛函之一。设 $M$ 紧致无边界，定义总标量曲率
\[
\mathcal S[g]
:=\int_M\mathcal R_g\,\operatorname{vol}_g.
\]
对任意对称二张量变分 $h=\dot g$，利用\eqnrefs{eq:math-dg-scalar-curvature-variation}和\eqnrefs{eq:math-dg-volume-variation}，有
\[
\begin{aligned}
\frac{\dd}{\dd t}\mathcal S[g(t)]
={}&\int_M
\left[
-\langle h,\operatorname{Ric}\rangle
+\nabla^i\nabla^jh_{ij}
-\Delta_{\mathrm{LB}}(\operatorname{tr}_gh)
\right]\operatorname{vol}_g\\
&+\frac12\int_M
\mathcal R\operatorname{tr}_g(h)\,\operatorname{vol}_g.
\end{aligned}
\]
闭流形上两个二阶导数项都是散度，积分为零，因此
\begin{equation*}\boxed{
\delta\mathcal S_g[h]
=-\int_M
\left\langle
\operatorname{Ric}-\frac12\mathcal R g,
h
\right\rangle_g\operatorname{vol}_g
}.\end{equation*}
括号中的
\[
G:=\operatorname{Ric}-\frac12\mathcal R g
\]
就是 Einstein 张量。它并不是人为组合出来的二阶张量，而是总标量曲率对度量的 $L^2$ 变分梯度（差一个符号）。

若只允许固定总体积的变分，需要加入约束
\[
\operatorname{Vol}(M,g)
=\int_M\operatorname{vol}_g=\text{const.}
\]
并用 Lagrange 乘子。临界方程变成
\[
\operatorname{Ric}-\frac12\mathcal R g
+\lambda g=0.
\]
取迹可解出 $\lambda$ 与 $\mathcal R$ 的关系，最终等价于
\begin{equation*}\boxed{
\operatorname{Ric}=cg
}\end{equation*}
其中 $c$ 为常数；这就是 Einstein 度量。于是 Einstein 条件的变分来源先于任何特定坐标公式。

Hodge 星同时依赖逆度量和体积，因此它的变化不是只乘一个标量。令 $\mathcal V:TM\to TM$ 为把 $\dot g$ 升一个指标所得的自伴 自同态，即由
\begin{equation*}g(\mathcal V X,Y)=\dot g(X,Y)
\end{equation*}
定义的自伴 自同态。对 $k$-形式 $\alpha$，定义
\[
(\mathcal V^{[k]}\alpha)(X_1,\ldots,X_k)
:=\sum_{r=1}^k
\alpha(X_1,\ldots,\mathcal V X_r,\ldots,X_k).
\]
若 $\alpha$ 本身不显含时间，则
\begin{equation}\label{eq:math-dg-hodge-star-variation}
\boxed{
\partial_t(*\alpha)
=*
\left[
\frac12\operatorname{tr}_g(\dot g)\,\alpha
-\mathcal V^{[k]}\alpha
\right].}
\end{equation}
若 $\alpha=\alpha(t)$，右端再加 $*\partial_t\alpha$。

\subsection{叶内体积与法向输运}

对随参数变化的超曲面，前面的 Lie 输运公式可以在适配坐标中逐项落实。

令叶片诱导度规为 $h_{ab}$，单位法向为 $n$，约定 $K_{ab}=h_a{}^\mu h_b{}^\nu\nabla_\mu n_\nu$，$K=h^{ab}K_{ab}$。它与嵌入章法向分量 $B_{ab}$ 的关系是 $K_{ab}=-\langle B_{ab},n\rangle$。

只要能够在局部用一族超曲面描述空间，场的微分方程就会自然分成沿法向的变化与沿超曲面的变化。此时需要进一步考察法向、诱导度规和体积元之间的关系。

定义第二基本形式 $K_{ab}=\frac12(\mathcal L_n h)_{ab}$，并约定 $K=h^{ab}K_{ab}$。直接对 \cref{eq:ghs-foliation-adm} 求 Lie 导数可得
\begin{equation}\label{eq:ghs-extrinsic-curvature}
 K_{ab}=\frac1{2N}\bigl(\partial_t h_{ab}-D_aN_b-D_bN_a\bigr),
\end{equation}
其中 $D$ 是 $h$ 的 Levi--Civita 联络。这既是几何恒等式，又说明“叶的形变”不仅是坐标面积变化，还包含切向重标记。

\begin{proposition}[法向体积变化]
设 $X=X^n n+X^a\partial_a$，则
\[
 \nabla_\mu X^\mu=\nabla_nX^n+K X^n+D_aX^a
 - X^\mu_{\parallel}\,n^\nu\nabla_\nu n_\mu.
\]
特别地，$\nabla_\mu n^\mu=K$，故在切向位移为零、法向进动系数为一的 规范下 $\partial_t\sqrt h=K\sqrt h$。
\end{proposition}
\begin{proof}
由 $n_\mu n^\mu=1$ 得 $n^\mu\nabla_\nu n_\mu=0$。将单位切向投影 $h^\mu{}_{\nu}=\delta^\mu{}_\nu-n^\mu n_\nu$ 插入 $\nabla_\mu n^\mu$，得到 $h^{\mu\nu}\nabla_\mu n_\nu=K$。利用 $\partial_t\log\sqrt h=\frac12h^{ab}\partial_t h_{ab}$，再代入 \cref{eq:ghs-extrinsic-curvature}，可得一般式 $\partial_t\sqrt h=\sqrt h(NK+D_aN^a)$。令 $a_\mu=n^\nu\nabla_\nu n_\mu$。由于 $n_\mu X^\mu_\parallel=0$，沿 $n$ 求导给出
\[
 n^\mu n^\nu\nabla_\mu(X_\parallel)_\nu=-a_\nu X^\nu_\parallel,
 \qquad \nabla_\mu X^\mu_\parallel=D_aX^a-a_aX^a.
\]
而 $\nabla_\mu(X^nn^\mu)=\nabla_nX^n+KX^n$。把两式相加就得到散度分解；此处 $D_aX^a$ 是叶内散度，与环境全散度不同。
\end{proof}

对于切向位移为零 的具体情形 $g=\dd t^2+h_{ab}(t,y)\dd y^a\dd y^b$，直接从
$\Delta_g\varphi=(\sqrt g)^{-1}\partial_\mu(\sqrt g\,g^{\mu\nu}\partial_\nu\varphi)$ 算出
\begin{equation}\label{eq:ghs-laplace-evolving-slices}
 \Delta_g\varphi=\partial_t^2\varphi+K\partial_t\varphi+\Delta_{h(t)}\varphi.
\end{equation}
第二项 $K\partial_t\varphi$ 正是叶的体积膨胀对波动传播的影响。在球面上，它退化为熟悉的 $(d-1)r^{-1}\partial_r$；但该式在一般弯曲叶上依然成立。


\section{纯微分同胚变化与 Ricci flow}

若度量变化只是由一族微分同胚产生，写
\[
g_t=\Psi_t^*g_0,
\qquad
W_t:=\partial_t\Psi_t\circ\Psi_t^{-1}.
\]
这里 $W_t$ 是作用在目标流形上的 Eulerian 生成元。由拉回微分公式，
\[
\partial_tg_t
=\Psi_t^*(\mathcal L_{W_t}g_0).
\]
利用 Lie 导数对微分同胚的自然性，若把生成元搬回参考流形并记
\[
U_t:=\Psi_t^*W_t,
\]
则同一式也可写成
\begin{equation*}\boxed{\partial_tg_t=\mathcal L_{U_t}g_t.}
\end{equation*}
由 Killing 方程推导中已经得到的恒等式，Levi--Civita 联络下
\begin{equation*}(\mathcal L_Ug)_{ij}
=\nabla_iU_j+\nabla_jU_i.
\end{equation*}
这种 $\dot g$ 只是在重新标记点，几何不变量随拉回一起变化，并不代表新的内在几何。这里必须区分 $W_t$ 与 $U_t=\Psi_t^*W_t$；除非额外满足相应不变性，不能把 $\Psi_t^*(\mathcal L_{W_t}g_0)$ 直接改写成 $\mathcal L_{W_t}g_t$。

真正的几何流则指定一个不纯粹来自 Lie 导数的 $\dot g$。最重要的例子之一是 Ricci 流
\begin{equation*}\partial_tg=-2\operatorname{Ric}(g).
\end{equation*}
它只是一般变分公式中的特殊代入。由\eqnrefs{eq:math-dg-volume-variation}与标量曲率定义，立刻得到
\begin{equation}\label{eq:math-dg-ricci-flow-volume}
\partial_t\operatorname{vol}_g
=-\mathcal R\,\operatorname{vol}_g.
\end{equation}
由\eqnrefs{eq:math-dg-connection-variation}则得到 Levi--Civita 联络的变化由 $\nabla\operatorname{Ric}$ 控制。将 $\dot g=-2\operatorname{Ric}$ 代入\eqnrefs{eq:math-dg-scalar-curvature-variation}，先得到
\[
\partial_t\mathcal R
=2|\operatorname{Ric}|^2
-2\nabla^i\nabla^j\operatorname{Ric}_{ij}
+2\Delta_{\mathrm{LB}}\mathcal R.
\]
再用 收缩 Bianchi 恒等式
\[
\nabla^i\operatorname{Ric}_{ij}=\frac12\nabla_j\mathcal R,
\]
得到
\begin{equation}\label{eq:math-dg-ricci-flow-scalar-curvature}
\boxed{
\partial_t\mathcal R
=\Delta_{\mathrm{LB}}\mathcal R+2|\operatorname{Ric}|^2.}
\end{equation}
因此 Ricci 流 下标量曲率的反应--扩散结构不是另行记忆的公式，而是一般度量变分的直接特例。

Ricci 流 仍具有微分同胚不变性，因此其坐标主符号带有 规范 退化。选定固定背景度量 $\bar g$，定义 DeTurck 向量场
\begin{equation*}W^k
=g^{ij}\left(\Gamma^k{}_{ij}(g)-\bar\Gamma^k{}_{ij}(\bar g)\right).
\end{equation*}
考虑 Ricci--DeTurck 方程
\begin{equation*}\boxed{
\partial_tg=-2\operatorname{Ric}(g)+\mathcal L_Wg.}
\end{equation*}
在局部坐标中，$-2\operatorname{Ric}$ 的非严格抛物部分恰被 $\mathcal L_Wg$ 的最高阶项抵消，留下
\[
\partial_tg_{ij}
=g^{ab}\partial_a\partial_bg_{ij}
+\text{至多一阶导数的非线性项}.
\]
所以固定 $\bar g$ 后，Ricci--DeTurck 系统具有严格抛物主部。求解它以后，再沿由 $-W$ 生成的时间依赖微分同胚拉回，就恢复 Ricci 流。DeTurck trick 因而只是固定微分同胚 规范，并没有改变内在几何内容。

Ricci 流 会改变总体积。若希望比较固定体积下的形状演化，可定义平均标量曲率
\[
\bar{\mathcal R}(t)
:=\frac{1}{\operatorname{Vol}(M,g(t))}
\int_M\mathcal R\,\operatorname{vol}_g
\]
并考虑 normalized Ricci 流
\begin{equation*}\boxed{
\partial_tg
=-2\operatorname{Ric}+\frac{2}{n}\bar{\mathcal R}\,g
}.\end{equation*}
由体积变分公式，
\[
\frac{\dd}{\dd t}\operatorname{Vol}(M,g)
=\int_M
\left(-\mathcal R+\bar{\mathcal R}\right)\operatorname{vol}_g=0.
\]
因此归一化项不是随意加入的尺度修正，而是精确消去全局体积漂移的 Lagrange-乘子 型项。

对标量曲率，代入一般变分公式得到
\begin{equation*}\partial_t\mathcal R
=\Delta_{\mathrm{LB}}\mathcal R
+2|\operatorname{Ric}|^2
-\frac{2}{n}\bar{\mathcal R}\,\mathcal R.\end{equation*}
若在某个时间区间内曲率有下界或上界，这类反应--扩散方程可以直接配合 maximum 原理。对非归一化 Ricci 流，若 $\mathcal R$ 在时刻 $t$ 取得空间最小值，则在该点 $\Delta_{\mathrm{LB}}\mathcal R\ge0$，并由
\[
|\operatorname{Ric}|^2\ge\frac1n\mathcal R^2
\]
得到
\[
\partial_t\mathcal R_{\min}
\ge\frac{2}{n}\mathcal R_{\min}^2
\]
的 barrier 意义估计。它说明 标量曲率 的负下界不会任意快向 $-\infty$ 漂移，而正曲率又可能在有限时间内因 reaction 项 增强。这类 maximum-原理 推导是几何流从张量 PDE 获得几何结论的基本方式。

标量曲率只记录曲率张量的一次收缩。若要研究 Ricci 流 是否形成奇性，必须继续跟踪 $\operatorname{Ric}$ 和完整 $\operatorname{Rm}$。把 $h=-2\operatorname{Ric}$ 代入一般 Ricci 变分式，并利用两次 收缩 Bianchi 恒等式整理，可得 $(0,2)$ 型 Ricci 张量 的精确演化
\begin{equation*}\boxed{
\partial_t\operatorname{Ric}_{ij}
=\Delta_{\mathrm{LB}}\operatorname{Ric}_{ij}
-2R_{ikj\ell}\operatorname{Ric}^{k\ell}
-2\operatorname{Ric}_i{}^k\operatorname{Ric}_{kj}.
}\end{equation*}
对它取迹时要同时考虑逆度量的变化
\[
\partial_tg^{ij}=2\operatorname{Ric}^{ij},
\]
这里 $R_{ikj\ell}=g(R(\partial_i,\partial_k)\partial_j,\partial_\ell)$，因此 $g^{ij}R_{ikj\ell}=-\operatorname{Ric}_{k\ell}$。
对 Ricci 演化式收缩时，两个二次项的迹互相抵消；再加上逆度量变化贡献的 $2|\operatorname{Ric}|^2$，便恢复\eqnrefs{eq:math-dg-ricci-flow-scalar-curvature}。
所以求一个随时间变化的张量的迹时，必须同时改变张量和收缩它的逆度量。

完整 Riemann 张量 的分量公式较长，但其结构非常稳定：
\begin{equation*}\boxed{
(\partial_t-\Delta_{\mathrm{LB}})\operatorname{Rm}
=\operatorname{Rm}*\operatorname{Rm},
}\end{equation*}
其中 $*$ 表示由度量收缩得到的有限个二次曲率组合。相应地，曲率范数满足
\begin{equation*}(\partial_t-\Delta_{\mathrm{LB}})|\operatorname{Rm}|^2
=-2|\nabla\operatorname{Rm}|^2+\operatorname{Rm}*\operatorname{Rm}*\operatorname{Rm}.\end{equation*}
真正重要的是右边只有“耗散梯度项 + 三次反应项”。这与半线性抛物 PDE 的 爆破 结构完全一致：扩散试图平滑曲率，非线性曲率自相互作用则可能放大高曲率区。

Ricci 流 具有精确抛物尺度。若 $g(t)$ 是解，则对任意 $\lambda>0$，
\begin{equation*}\widetilde g(t)
:=\lambda g(t/\lambda)\end{equation*}
仍是 Ricci 流。长度尺度变为 $\sqrt\lambda$，时间尺度变为 $\lambda$，而
\[
|\widetilde{\operatorname{Rm}}|_{\widetilde g}
=\lambda^{-1}|\operatorname{Rm}|_g.
\]
所以在高曲率点附近选择与曲率大小相匹配的抛物缩放，就能把该点曲率归一到 $O(1)$。几何流奇性分析中的 切 流、ancient 解 与 爆破 极限 都来自这一尺度律，而不是额外的技巧。

对闭流形，Ricci--DeTurck 的严格抛物性给出短时存在唯一性。更进一步，若最大存在区间为 $[0,T)$ 且 $T<\infty$，则 Hamilton 扩张 criterion 说明
\begin{equation*}\boxed{
\limsup_{t\uparrow T}
\sup_{x\in M}|\operatorname{Rm}(x,t)|=\infty.
}\end{equation*}
换言之，有限时间失效不能只由坐标或微分同胚 规范 造成；必定有真正的内在曲率 爆破。这个结论与非线性 PDE 章节中的 continuation criterion 具有相同逻辑：只要控制决定局部适定性的最高阶几何量，就能把解延拓过当前时间。

二维情形把这些结构压缩得最透明。若 $n=2$，则
\[
\operatorname{Ric}=\frac12\mathcal R g,
\]
故 Ricci 流 变成
\begin{equation*}\partial_tg=-\mathcal R g,\end{equation*}
而标量曲率满足
\begin{equation*}\boxed{
\partial_t\mathcal R
=\Delta_{\mathrm{LB}}\mathcal R+\mathcal R^2.
}\end{equation*}
这说明二维 Ricci 流 本质上是一个带几何耦合的反应--扩散方程；高维中新出现的困难正来自 Weyl 曲率 与 Ricci 张量 不再由一个标量决定。
Ricci 流 的 平稳 point 不能只写成 $\operatorname{Ric}=0$，因为几何流允许同时发生整体缩放与微分同胚重标记。把这两个“平凡运动”都除掉后，自相似解由 Ricci 孤子 方程刻画：存在向量场 $X$ 与常数 $\lambda$ 使
\begin{equation}
\boxed{
\operatorname{Ric}+\frac12\mathcal L_Xg=\lambda g.
}
\label{eq:math-dg-ricci-soliton}
\end{equation}
$\lambda>0,=0,<0$ 时分别称 shrinking、steady、expanding 孤子。若
\[
X=\nabla f,
\]
则
\[
\frac12\mathcal L_{\nabla f}g=\nabla^2f,
\]
于是得到 梯度 Ricci 孤子
\begin{equation}
\boxed{
\operatorname{Ric}+\nabla^2f=\lambda g.
}
\label{eq:math-dg-gradient-ricci-soliton}
\end{equation}

按 $\lambda$ 的符号，孤子 分为三类：Shrinker（$\lambda>0$，流在 $t\to1/(2\lambda)$ 时收缩到一点，是 Ricci 流奇点的 爆破 模型，典型如 Gauss 收缩子 $g_{\mathrm{euc}}$、$X=-\tfrac12\nabla|x|^2$）、Steady（$\lambda=0$，流仅由微分同胚驱动，体积不变，典型如 cigar 孤子 $g=(\dd x^2+\dd y^2)/(1+x^2+y^2)$）、Expander（$\lambda<0$，流对所有 $t<0$ 存在，从 $t\to-\infty$ 的一点展开）。当 $X=\nabla f$ 为梯度场时，孤子 方程简化为 $\operatorname{Ric}+\nabla^2f=\lambda g$；取迹得 $\mathcal R+\Delta f=n\lambda$，对散度并利用 收缩 Bianchi 恒等式 $\nabla^i\operatorname{Ric}_{ij}=\tfrac12\nabla_j\mathcal R$ 得首积分 $\mathcal R+|\nabla f|^2-2\lambda f=C$，即 Perelman $\mathcal W$ 泛函的临界点方程。$\mathcal W[g,f,\tau]=\int_M[\tau(\mathcal R+|\nabla f|^2)-f-n]\ee^{-f}/(4\pi\tau)^{n/2}\operatorname{vol}_g$ 在 Ricci 流耦合下单调不增，临界点恰为 梯度 收缩子，故 收缩子 是 Ricci 流的'熵极大点'。例：$\RR^n$ 上 $g_{\mathrm{euc}}$、$f=|x|^2/4\tau$、$\lambda=1/2\tau$ 直接验证满足 收缩子 方程（$\operatorname{Ric}=0$、$\nabla^2f=g/2\tau$），对应热方程基本解；二维 cigar 孤子 $g=(\dd x^2+\dd y^2)/(1+x^2+y^2)$、$X=-2(x\partial_x+y\partial_y)$、$\lambda=0$ 是 稳态孤子，曲率 $\mathcal R=4/(1+r^2)$ 在无穷远趋于零。Bryant 孤子 是三维 Ricci 流中非紧 收缩子 的典型例子，渐近圆柱形；广义相对论中梯度 Ricci 孤子 方程与静态 Einstein 方程通过 共形 变换相关，例如 Schwarzschild 几何在适当坐标下满足 孤子 型方程。

自相似解的流参数化由此给出：$\phi_t$ 由 $Y_t=X/c(t)$ 生成（\eqnrefs{eq:math-dg-soliton-selfsimilar-flow}）。

对 梯度孤子，取迹得到
\begin{equation*}\mathcal R+\Delta_{\mathrm{LB}}f=n\lambda.\end{equation*}
对\eqnrefs{eq:math-dg-gradient-ricci-soliton} 取散度，并使用 收缩 Bianchi 恒等式与 Hessian 的交换公式，可得
\[
\dd\mathcal R=2\operatorname{Ric}(\nabla f,\cdot).
\]
于是
\begin{equation*}\boxed{
\mathcal R+|\nabla f|^2-2\lambda f=C
}\end{equation*}
在每个连通分支上为常数。这个 first 积分 在研究 收缩子 的加权体积、势函数增长与 爆破 极限 时非常重要。

自相似结构还可以由一个单调泛函识别。对闭流形定义 Perelman $\mathcal F$-泛函
\begin{equation*}\mathcal F(g,f)
:=\int_M
\left(\mathcal R+|\nabla f|^2\right)
\ee^{-f}\operatorname{vol}_g,
\qquad
\int_M\ee^{-f}\operatorname{vol}_g=1.\end{equation*}
若 $(g,f)$ 满足耦合系统
\begin{equation}
\partial_tg=-2\operatorname{Ric},
\qquad
\partial_tf
=-\Delta_{\mathrm{LB}}f+|\nabla f|^2-\mathcal R,
\label{eq:math-dg-perelman-coupled-flow}
\end{equation}
则加权测度总质量保持，并且
\begin{equation}
\boxed{
\frac{\dd}{\dd t}\mathcal F(g,f)
=2\int_M
|\operatorname{Ric}+\nabla^2f|^2
\ee^{-f}\operatorname{vol}_g
\ge0.
}
\label{eq:math-dg-perelman-f-monotonicity}
\end{equation}
因此 Ricci 流 并不是一个没有 Lyapunov 结构的任意非线性 PDE；在与适当 backward 热-型 势函数耦合后，它具有严格的熵型单调量。

先验证约束确实保持。由\eqnrefs{eq:math-dg-ricci-flow-volume} 与\eqnrefs{eq:math-dg-perelman-coupled-flow}，
\begin{align*}
\frac{\dd}{\dd t}\int_M\ee^{-f}\operatorname{vol}_g
&=\int_M(-\partial_tf-\mathcal R)\ee^{-f}\operatorname{vol}_g\\
&=\int_M
(\Delta_{\mathrm{LB}}f-|\nabla f|^2)
\ee^{-f}\operatorname{vol}_g=0,
\end{align*}
其中最后一步用分部积分
\[
\int_M\ee^{-f}\Delta_{\mathrm{LB}}f\,\operatorname{vol}_g
=\int_M\ee^{-f}|\nabla f|^2\operatorname{vol}_g.
\]
而\eqnrefs{eq:math-dg-perelman-f-monotonicity} 则来自标量曲率、梯度能与加权体积三项变分的精确抵消；剩余项完成平方为
\[
2|\operatorname{Ric}+\nabla^2f|^2.
\]
这与前面“总标量曲率的一阶变分给 Einstein 张量”形成对应：这里不是寻找静态 Einstein 临界点，而是在加权几何中找到 Ricci 流 的梯度型结构。

对 收缩子 的完整刻画还需要加入尺度归一化后的 $\mathcal W$-熵；本章先保留到 $\mathcal F$ 这一层，因为它已足以说明三个核心事实：微分同胚 规范 由 DeTurck 处理，自相似 爆破 模型由 Ricci 孤子 描述，而熵型单调量为排除某些回返行为和分析奇性极限提供全局方向。

\begin{exbox}[闭形式的冻结输运]
设 $\alpha_t\in\Omega^k(M)$ 满足
\[
\partial_t\alpha_t+\mathcal L_{V_t}\alpha_t=0.
\]
由\eqnrefs{eq:math-dg-geometric-transport-solution}，
\[
\Phi_{t,s}^*\alpha_t=\alpha_s.
\]
若 $\dd\alpha_s=0$，则
\[
\Phi_{t,s}^*(\dd\alpha_t)
=\dd(\Phi_{t,s}^*\alpha_t)
=\dd\alpha_s=0.
\]
微分同胚拉回可逆，所以 $\dd\alpha_t=0$。因此初始闭形式在随流输运后仍闭；理想流体和磁流体中的“冻结形式”正是按这一拉回关系演化。
\end{exbox}
输运恒等式区分坐标拖曳与度量本身的变化。以下将两者代入体积元和曲率的变分，明确比较对象时需要固定哪些结构。

\section{度量变化与坐标拖曳}

沿一个向量场移动观察点，会改变度量的坐标表达式，即使曲面本身没有改变。要区分这种搬运与真正的几何演化，先把不同时间的张量拉回到同一点比较。

\begin{definition}[Lie 导数与度量变分]
设 \(X\) 生成局部流 \(\Phi_t\)。张量场 \(T\) 的 Lie 导数为
\(\mathcal L_XT=\left.\frac{\dd}{\dd t}\right|_{t=0}\Phi_t^*T\)。
一族度量 \(g(t)\) 的变分张量为 \(h=\partial_tg\)，即每时每点上的对称双线性型。\(\mathcal L_Xg\) 是由微分同胚流产生的一种特殊变分。
\end{definition}

例如平面上的平移流不改变 欧氏 度量，故 \(\mathcal L_{\partial_x}g=0\)；伸缩流 \(\Phi_t(x)=\ee^t x\) 则满足 \(\Phi_t^*g=\ee^{2t}g\)，从而 \(\mathcal L_Xg=2g\)。伸缩前后的度量分量不同，却仍彼此等距，因为它们由 \(\Phi_t\) 联系。若时间演化不能写成拉回，才可能产生不能通过重新标记点消去的变化。

这里用到三个须先固定符号的缩并。对 Levi--Civita 曲率
\(R(X,Y)Z\)，定义 Ricci 张量为
\(\operatorname{Ric}(Y,Z)=
\operatorname{tr}[X\mapsto R(X,Y)Z]\)，
即把第一个槽位 \(X\) 与输出向量的求导指标缩并，
保留的第二、三槽位 \(Y,Z\) 成为 \(\operatorname{Ric}\)
的两个自变量；若改用相反号差的曲率约定，
\(\operatorname{Ric}\) 将整体变号。
再定义标量曲率
\(R_g=\operatorname{tr}_g\operatorname{Ric}
=g^{ij}\operatorname{Ric}_{ij}\)。
对函数取 \(\Delta_g f=\nabla^i\nabla_i f\)，
所以在 欧氏 坐标中 \(\Delta_g=\sum_i\partial_i^2\)，
与上一章非负的 \(\Delta_{\rm dR}\) 在函数上相差一个负号。
在二维单位球面，前章算出的 \(K=1\) 给
\(\operatorname{Ric}=g\)、\(R_g=2\)，
这也固定了下面变分式的曲率号差。

\begin{proposition}[自然标量随拉回而变]
设 \(\Phi:M\to M\) 为微分同胚，\(I(g)\) 是由 \(g\)、其 Levi--Civita 联络、曲率及张量缩并自然构成的标量，例如标量曲率。则 \(I(\Phi^*g)=\Phi^*I(g)\)。特别地，若 \(g(t)=\Phi_t^*g_0\)，则 \(I(g(t))(p)=I(g_0)(\Phi_t(p))\)。
\end{proposition}
\begin{proof}
\(\Phi\) 把 \((M,\Phi^*g)\) 等距映到 \((M,g)\)。拉回的联络仍无挠且与拉回度量相容；Levi--Civita 唯一性说明它正是 \(\Phi^*g\) 的联络。曲率由联络交换子定义，故随之拉回；张量缩并与拉回可交换。对所得标量求值即得结论。
\end{proof}

对任意变分 \(h=\partial_tg\)，体积元满足
\[
\partial_t(\dd V_g)=\tfrac12\operatorname{tr}_g(h)\,\dd V_g.
\]
因为 \(\dd V_g=\sqrt{\det(g_{ij})}\,\dd x^1\cdots \dd x^n\)，而矩阵恒等式
\(\partial_t\log\det g=\operatorname{tr}(g^{-1}h)\) 来自行列式一阶展开。
曲率的变化需要再多一步，因为联络本身也随度量改变。对
\(g(t)\) 在固定 \(t\) 处求导，记 \(h_{ij}=\partial_tg_{ij}\)、
\(h^{ij}=g^{ik}g^{j\ell}h_{k\ell}\)。从
\(g^{ik}g_{kj}=\delta^i_j\) 求导，得到
\(\partial_tg^{ij}=-h^{ij}\)。再把
\(\partial_t(\nabla_i g_{jk})=0\) 与联络无挠条件放进上一章的
Koszul 公式，三式相加减，得到
\[
\partial_t\Gamma^k{}_{ij}
=\tfrac12g^{k\ell}
\bigl(\nabla_i h_{j\ell}+\nabla_j h_{i\ell}-\nabla_\ell h_{ij}\bigr).
\]
这里右端全是同一时刻 \(g\) 的协变导数；\(\partial_t\Gamma\) 是张量，虽然单独的 \(\Gamma\) 不是张量，所以这个表达在换坐标后仍一致。

沿用 \(R(X,Y)Z=\nabla_X\nabla_YZ-\nabla_Y\nabla_XZ-\nabla_{[X,Y]}Z\) 的约定，变分曲率中两个 \(\Gamma\Gamma\) 项与偏导项恰组合为协变导数，因而
\(\partial_t\operatorname{Ric}_{ij}
=\nabla_k(\partial_t\Gamma^k{}_{ij})
-\nabla_j(\partial_t\Gamma^k{}_{ik})\)。把刚才的联络变分代入，并用
\(\partial_tR=(\partial_tg^{ij})\operatorname{Ric}_{ij}
+g^{ij}\partial_t\operatorname{Ric}_{ij}\)。具体地，先把联络变分按
\(\partial_t\Gamma^k{}_{ij}
=\tfrac12g^{k\ell}(\nabla_i h_{j\ell}
+\nabla_j h_{i\ell}-\nabla_\ell h_{ij})\)
逐项缩并。对第二项缩并，哑指标在两项间改名后相消：
\[
\partial_t\Gamma^k{}_{ik}
=\tfrac12g^{k\ell}\nabla_i h_{k\ell}
+\tfrac12\bigl(\nabla^\ell h_{i\ell}-\nabla^k h_{ik}\bigr)
=\tfrac12\nabla_i(\operatorname{tr}_g h),
\]
其中括号内第二、三项由 \(h\) 的对称性、经改名
\(\ell\leftrightarrow k\) 后抵消，首项由度量相容性收成
\(\operatorname{tr}_g h=g^{k\ell}h_{k\ell}\) 的导数。同理
\[
g^{ij}\partial_t\Gamma^k{}_{ij}
=\tfrac12g^{k\ell}\bigl(2\nabla^j h_{j\ell}
-g^{ij}\nabla_\ell h_{ij}\bigr)
=\nabla_i h^{ik}-\tfrac12\nabla^k(\operatorname{tr}_g h),
\]
这里前两项因 \(h_{ij}=h_{ji}\) 相等（把哑指标
\(i\leftrightarrow j\) 改名即见）。对两式分别再求
\(\nabla_k\) 与 \(\nabla^i\)，得到
\begin{equation}
\label{eq:math-scalar-curvature-variation}
\partial_tR
=-\langle\operatorname{Ric},h\rangle_g
+\nabla^i\nabla^j h_{ij}-\Delta_g(\operatorname{tr}_g h).
\end{equation}
最后两项是散度型；它们在有边界的区域上积分时留下边界贡献，不能不加条件地丢掉。作为量纲与符号检查，若 \(g(t)=a(t)^2g_0\) 且 \(a\) 只依赖时间，便有 \(h=2a'/a\,g\) 与 \(\nabla h=0\)。式 \cref{eq:math-scalar-curvature-variation} 化为
\(\partial_tR=-2a'/a\,R\)，正与 \(R(g(t))=a(t)^{-2}R(g_0)\) 一致。
若取 Ricci 流 \(h=-2\operatorname{Ric}(g)\)，便有
\(\partial_t \dd V_g=-R_g\dd V_g\)。若演化方程右端另加一项
\(\mathcal L_Xg\)，由上面关于拉回的命题，那一项只对应把点重新命名，并不改变任何曲率。

半径随时间变化的球面能把这些式子算到头。令
\(g(t)=a(t)^2g_{S^2(1)}\)，其中单位球面度量
\(g_{S^2(1)}=\dd\vartheta^2+\sin^2\vartheta\,\dd\varphi^2\)。
前面的曲率计算给
\(K(t)=a(t)^{-2}\)、\(R_g=2K=2/a(t)^2\)；
在二维 \(\operatorname{Ric}=Kg=g_{S^2(1)}\)。
把这些量代入 Ricci 流方程，逐项比较
\[
\partial_tg=\frac{\dd}{\dd t}a(t)^2\,g_{S^2(1)}
=-2g_{S^2(1)}
\]
的系数，得到
\(a(t)^2=a(0)^2-2t\)，只在 \(0\le t<a(0)^2/2\)
保持正定。球面面积为 \(4\pi a(t)^2\)，所以面积变化率为
\(-8\pi\)；另一方面
\(-\int_{S^2}R_g\,\dd V_g
=-(2/a^2)(4\pi a^2)=-8\pi\)，
正好验证体积元公式。到 \(t=a(0)^2/2\) 时半径缩为零，
此处光滑度量不再存在，方程不能继续按同一光滑解解释。



\par
\begin{workedexample}
在二维球面上令 \(g(t)=a(t)^2g_{S^2(1)}\)，计算
\(\partial_t \dd V_g\) 与 \(R_g\)；再求它满足 Ricci 流时 \(a(t)^2\) 的变化。
\textbf{解.} \(g(t)=a^2g_0\) 给 \(\dd V_g=a^2\dd V_0\)，因而
\[
\partial_t\dd V_g=2\frac{a'}a\dd V_g.
\]
二维曲率在度量整体放大 \(a^2\) 时按 \(a^{-2}\) 缩放，故 \(R_g=2/a^2\)。又二维有 \(\operatorname{Ric}(g)=g_0\)，Ricci 流 \(\partial_tg=-2\operatorname{Ric}(g)\) 化为 \((a^2)'g_0=-2g_0\)，即 \(a(t)^2=a(0)^2-2t\)（至曲率奇点时间）。
\end{workedexample}
\par



\section{可积条件与几何演化的公式联系}
Frobenius 条件描述固定时刻的可积性，输运公式描述随时间变化的张量。以下取光滑常秩分布；余维一的定义形式 $\alpha$ 在讨论邻域内处处非零。度量变分中 $\dot g^{ij}=g^{ia}g^{jb}\partial_tg_{ab}$ 表示先变分再升指标。
\begin{note}[公式之间的联系]
Frobenius 的向量场与余切判据分别为
\[
[\Gamma(\mathcal D),\Gamma(\mathcal D)]\subset\Gamma(\mathcal D),
\qquad
\dd\alpha^a=\eta^a{}_b\wedge\alpha^b.
\]
余维一时可写成
\[
\alpha\wedge\dd\alpha=0.
\]
时间依赖流与度量变分的核心公式为
\begin{align*}
\frac{\dd}{\dd t}(\Phi_{t,s}^*T_t)
&=\Phi_{t,s}^*(\partial_tT_t+\mathcal L_{V_t}T_t),\\
\partial_tg^{ij}&=-\dot g^{ij},\\
\partial_t\operatorname{vol}_g&=\frac12\operatorname{tr}_g(\dot g)\operatorname{vol}_g,\\
\partial_t\Gamma^k{}_{ij}
&=\frac12g^{k\ell}(\nabla_i \dot g_{j\ell}+\nabla_j \dot g_{i\ell}-\nabla_\ell \dot g_{ij}).
\end{align*}
完整条件见\eqnrefs{eq:math-dg-transport-identity}、\eqnrefs{eq:math-dg-volume-variation}与\eqnrefs{eq:math-dg-connection-variation}。
\end{note}

\begin{proposition}[时间依赖张量的拉回微分]
设 $\Phi_{t,s}$ 为光滑时间依赖向量场 $V_t$ 的两参数局部流，$T_t$ 为光滑张量族。则\[\frac{\dd}{\dd t}(\Phi_{t,s}^*T_t)=\Phi_{t,s}^*(\partial_tT_t+\mathcal L_{V_t}T_t).\]
\end{proposition}

\begin{proof}
利用两参数流的组合律将 $\Phi_{t+\varepsilon,s}^*$ 分解为 $\Phi_{t,s}^*\circ\Phi_{t+\varepsilon,t}^*$，对短时流 $\Phi_{t+\varepsilon,t}^*$ 用 Lie 导数的一阶展开。

组合律分解。固定 $t$，取小量 $\varepsilon$。两参数流的组合律给出
\begin{equation*}
 \Phi_{t+\varepsilon,s}
 =\Phi_{t+\varepsilon,t}\circ\Phi_{t,s}.
\end{equation*}
因此拉回分解为
\begin{equation*}
 \Phi_{t+\varepsilon,s}^*T_{t+\varepsilon}
 =\Phi_{t,s}^*
 \left(\Phi_{t+\varepsilon,t}^*T_{t+\varepsilon}\right).
\end{equation*}

短时流展开。在时间 $t$ 附近，短时流 $\Phi_{t+\varepsilon,t}$ 由向量场 $V_t$ 生成，其拉回的一阶展开为
\begin{equation*}
 \Phi_{t+\varepsilon,t}^*
 =I+\varepsilon\,\mathcal L_{V_t}+o(\varepsilon),
\end{equation*}
而张量场的显式时间变化为
\begin{equation*}
 T_{t+\varepsilon}
 =T_t+\varepsilon\,\partial_tT_t+o(\varepsilon).
\end{equation*}

相乘并取差商。
\begin{align*}
 \Phi_{t+\varepsilon,t}^*T_{t+\varepsilon}
 &=\bigl(I+\varepsilon\,\mathcal L_{V_t}+o(\varepsilon)\bigr)
  \bigl(T_t+\varepsilon\,\partial_tT_t+o(\varepsilon)\bigr)\\
 &=T_t+\varepsilon
  \bigl(\partial_tT_t+\mathcal L_{V_t}T_t\bigr)+o(\varepsilon).
\end{align*}
因此
\begin{align*}
 \frac{\dd}{\dd t}\bigl(\Phi_{t,s}^*T_t\bigr)
 &=\lim_{\varepsilon\to0}\frac{\Phi_{t+\varepsilon,s}^*T_{t+\varepsilon}-\Phi_{t,s}^*T_t}{\varepsilon}\\
 &=\Phi_{t,s}^*
 \bigl(\partial_tT_t+\mathcal L_{V_t}T_t\bigr),
\end{align*}
即\eqnrefs{eq:math-dg-transport-identity}。
\end{proof}

\begin{proposition}[Levi--Civita 联络的度量变分]
设 $g_t$ 为光滑度量族，$h=\dot g$、$\mathcal A=\partial_t\nabla^{g_t}$。则 $\mathcal A$ 为张量，且\[\begin{aligned}2g(\mathcal A(X,Y),Z)={}&(\nabla_Xh)(Y,Z)+(\nabla_Yh)(X,Z)\\&-(\nabla_Zh)(X,Y).\end{aligned}\]
\end{proposition}

\begin{proof}
在每个固定 $t$ 上使用 Koszul 公式。向量场和 Lie 括号不随 $t$ 变化，所以对时间求导时只有度量与联络贡献。对
\[
2g(\nabla_XY,Z)
=Xg(Y,Z)+Yg(Z,X)-Zg(X,Y)
+g([X,Y],Z)-g([Y,Z],X)+g([Z,X],Y)
\]
求导，左边为
\[
2\dot g(\nabla_XY,Z)+2g(\mathcal A(X,Y),Z).
\]
右边把每个 $g$ 替换为 $\dot g$。再用当前 Levi--Civita 联络把普通方向导数拆成 $\nabla \dot g$ 与对向量槽位的联络作用；所有含 $\dot g(\nabla_XY,Z)$ 一类的项与左边对应项相消，得到
\begin{equation}\label{eq:math-dg-connection-variation}
\boxed{
\begin{aligned}
2g(\mathcal A(X,Y),Z)
={}&(\nabla_X\dot g)(Y,Z)
+(\nabla_Y\dot g)(X,Z)\\
&-(\nabla_Z\dot g)(X,Y).
\end{aligned}}
\end{equation}
坐标中即
\begin{equation}\label{eq:math-dg-christoffel-variation}
\boxed{
\partial_t\Gamma^k{}_{ij}
=\frac12g^{k\ell}
\left(
\nabla_i \dot g_{j\ell}
+\nabla_j \dot g_{i\ell}
-\nabla_\ell \dot g_{ij}
\right).}
\end{equation}

逐项消去的细节。对右边第一项 $Xg(Y,Z)$ 求导，先用 $\dot g$ 替换 $g$ 得 $\dot g(Y,Z)$ 的方向导数 $X\dot g(Y,Z)$，再用 Levi--Civita 联络展开
\begin{equation*}
X\dot g(Y,Z)
=(\nabla_X\dot g)(Y,Z)
+\dot g(\nabla_XY,Z)+\dot g(Y,\nabla_XZ).
\end{equation*}
对其他两项作类似展开。含 $\dot g(\nabla_XY,Z)$ 的项与左边 $2\dot g(\nabla_XY,Z)$ 部分相消：右边三项给出 $+1+1+0=2$ 倍 $\dot g(\nabla_XY,Z)$（前两项含 $\nabla_XY$，第三项含 $\nabla_ZX$ 不贡献），与左边相消。剩余项整理即得协变导数形式。

张量性的直接验证。要确认 $\mathcal A(X,Y)$ 是张量，需验证它对 $X,Y$ 都是函数线性的。对 $f\in C^\infty(M)$，
\begin{equation*}
\mathcal A(fX,Y)=\partial_t(\nabla_{fX}Y)=\partial_t(f\nabla_XY)=f\,\partial_t(\nabla_XY)=f\,\mathcal A(X,Y),
\end{equation*}
因 $f$ 与 $t$ 无关。对第二槽须显式保留联络的 Leibniz 项：
\[
\mathcal A(X,fY)=\partial_t\bigl(X(f)Y+f\nabla_XY\bigr)
=f\mathcal A(X,Y),
\]
因为 $X(f)Y$ 时间无关。这印证了"两个联络之差是张量"的一般事实。

Ricci 流 $\partial_tg=-2\operatorname{Ric}$ 给出 $\dot g=-2\operatorname{Ric}$，代入\eqnrefs{eq:math-dg-christoffel-variation} 得
\begin{equation*}
\partial_t\Gamma^k{}_{ij}
=-g^{k\ell}\left(
\nabla_i\operatorname{Ric}_{j\ell}
+\nabla_j\operatorname{Ric}_{i\ell}
-\nabla_\ell\operatorname{Ric}_{ij}
\right).
\end{equation*}

Hodge 星同时依赖 $g^{-1}$ 与体积形式 $\operatorname{vol}_g$，故其变分须合并\eqnrefs{eq:math-dg-inverse-metric-variation}与\eqnrefs{eq:math-dg-volume-variation}，下框给出完整公式。
\end{proof}

\begin{proposition}[曲率与标量曲率的变分]
对 $h=\dot g$、$\mathcal A=\dot\nabla$，\[\dot R(X,Y)Z=(\nabla_X\mathcal A)(Y,Z)-(\nabla_Y\mathcal A)(X,Z),\]\[\dot{\mathcal R}=-\langle h,\operatorname{Ric}\rangle_g+\nabla^i\nabla^jh_{ij}-\Delta_{\mathrm{LB}}\operatorname{tr}_gh.\]
\end{proposition}

\begin{proof}
从曲率定义
\[
R(X,Y)Z
=\nabla_X\nabla_YZ-\nabla_Y\nabla_XZ-\nabla_{[X,Y]}Z
\]
对时间求导。取 $X,Y,Z$ 为时间无关向量场，且 $\partial_t\nabla=\mathcal A$，得到
\[
\begin{aligned}
\partial_tR(X,Y)Z
={}&\mathcal A(X,\nabla_YZ)+\nabla_X\mathcal A(Y,Z)\\
&-\mathcal A(Y,\nabla_XZ)-\nabla_Y\mathcal A(X,Z)
-\mathcal A([X,Y],Z).
\end{aligned}
\]
另一方面，
\[
(\nabla_X\mathcal A)(Y,Z)
=\nabla_X(\mathcal A(Y,Z))
-\mathcal A(\nabla_XY,Z)
-\mathcal A(Y,\nabla_XZ).
\]
对 $X,Y$ 两项相减，并使用 Levi--Civita 联络无挠，
\[
\nabla_XY-\nabla_YX=[X,Y],
\]
便逐项恢复前一展开，从而得到\eqnrefs{eq:math-dg-curvature-variation-abstract}。在坐标基上取迹即得\eqnrefs{eq:math-dg-ricci-variation}。

最后对 $\mathcal R=g^{jk}\operatorname{Ric}_{jk}$ 求导：
\[
\partial_t\mathcal R
=(\partial_tg^{jk})\operatorname{Ric}_{jk}
+g^{jk}\partial_t\operatorname{Ric}_{jk}.
\]
第一项由\eqnrefs{eq:math-dg-inverse-metric-variation}给出
$-\langle\dot g,\operatorname{Ric}\rangle_g$。第二项把\eqnrefs{eq:math-dg-christoffel-variation}代入\eqnrefs{eq:math-dg-ricci-variation}，利用 $\nabla g=0$ 收缩指标，得到
\[
g^{jk}\partial_t\operatorname{Ric}_{jk}
=\nabla^i\nabla^j\dot g_{ij}
-\Delta_{\mathrm{LB}}(\operatorname{tr}_g\dot g),
\]
两项相加即为\eqnrefs{eq:math-dg-scalar-curvature-variation}。

标量曲率变分公式\eqnrefs{eq:math-dg-scalar-curvature-variation} 是几何流理论的基础：代入 Ricci 流 $\dot g=-2\operatorname{Ric}$ 得 $\partial_t\mathcal R=\Delta\mathcal R+2|\operatorname{Ric}|^2$，即标量曲率沿 Ricci 流 的反应--扩散方程。

共形变分 $\dot g=2u\,g$ 下，$\operatorname{tr}_g\dot g=2nu$、$\nabla^i\nabla^j\dot g_{ij}=2\Delta u$，代入得 $\partial_t\mathcal R=-2(n-1)\Delta_{\mathrm{LB}}u-2u\mathcal R$；$n=2$ 时为 $\partial_t\mathcal R=-2\Delta_{\mathrm{LB}}u-2u\mathcal R$，即 Gauss 曲率的共形演化。
\end{proof}

\begin{theorem}[固定体积的总标量曲率临界点]
设 $M^n$ 紧致、无边界，$n\ge3$。总标量曲率 $\mathcal S[g]=\int_M\mathcal R_g\operatorname{vol}_g$ 在固定总体积的度量空间上的临界点恰为 Einstein 度量 $\operatorname{Ric}=cg$。二维时该泛函由 Gauss--Bonnet 固定，不能由此推出常曲率。
\end{theorem}

\begin{proof}
对体积泛函
\[
\mathcal V[g]=\int_M\operatorname{vol}_g
\]
有
\[
\delta\mathcal V_g[h]
=\frac12\int_M\operatorname{tr}_g(h)\,\operatorname{vol}_g
=\frac12\int_M\langle g,h\rangle_g\operatorname{vol}_g.
\]
临界条件
\[
\delta(\mathcal S+2\lambda\mathcal V)[h]=0
\quad\forall h
\]
给出
\[
-\operatorname{Ric}+\frac12\mathcal R g+\lambda g=0.
\]
在 $n$ 维中取迹：
\[
-\mathcal R+\frac n2\mathcal R+n\lambda=0,
\]
$n\ge3$ 时可解出 $\mathcal R=-2n\lambda/(n-2)$，为常数。二维时这个取迹式只有 $\lambda=0$，并不能推出 $\mathcal R$ 恒定。代回原方程，得到
\[
\operatorname{Ric}=\left(\frac12\mathcal R+\lambda\right)g=:cg.
\]
反过来任何 Einstein 度量 满足带适当乘子的 Euler--Lagrange 方程。

体积变分的显式推导。度量变分 $g_t=g+th$ 下体积形式的变化为
\begin{equation*}
\operatorname{vol}_{g_t}
=\sqrt{\frac{\det(g+th)}{\det g}}\,\operatorname{vol}_g
=\sqrt{\det(I+tg^{-1}h)}\,\operatorname{vol}_g.
\end{equation*}
对 $t$ 在 $t=0$ 求导，用 $\det(I+A)$ 的一阶展开 $\det(I+A)=1+\operatorname{tr}A+O(\|A\|^2)$，
\begin{equation*}
\left.\frac{\dd}{\dd t}\operatorname{vol}_{g_t}\right|_{t=0}
=\frac12\operatorname{tr}(g^{-1}h)\,\operatorname{vol}_g
=\frac12\operatorname{tr}_g(h)\,\operatorname{vol}_g.
\end{equation*}
积分即得体积泛函的变分公式。$\operatorname{tr}_g(h)=g^{ij}h_{ij}$ 是关于 $g$ 的迹，与坐标无关。

常数 $c$ 的符号给出 Ricci 曲率的分类：$c>0$ 为正 Ricci（例如球面），$c=0$ 为 Ricci 平坦，$c<0$ 为负 Ricci（例如双曲空间）。Einstein 条件在维数大于三时不意味着所有截面曲率相同，不能据 $c$ 的符号直接分类全部截面曲率。Einstein 度量还是 Ricci 流 的缩放固定点（$g(t)=(1-2ct)g$），并在广义相对论中对应带宇宙学常数的真空 Einstein 方程。
\end{proof}

\begin{proposition}[Hodge 星的度量变分]
设 $g_t$ 为定向流形上的光滑度量族，$\mathcal V=g^{-1}\dot g$，$\alpha$ 为时间无关 $k$-形式。则\[\partial_t(*\alpha)=*\left[\tfrac12\operatorname{tr}_g\dot g\,\alpha-\mathcal V^{[k]}\alpha\right].\]特别地，$\dot g=2ug$ 时 $\partial_t(*\alpha)=(n-2k)u*\alpha$。
\end{proposition}

\begin{proof}
Hodge 星由度量通过内积和体积形式隐式定义。对度量变分，分别追踪内积变化和体积变化，再用 Hodge 星定义反读。

Hodge 星的隐式定义。对任意时间无关 $k$-形式 $\alpha,\beta$，Hodge 星由
\begin{equation*}
 \alpha\wedge *\beta
 =\langle\alpha,\beta\rangle_g\,\operatorname{vol}_g
\end{equation*}
唯一确定。两边对 $t$ 求导，$\alpha$ 时间无关故只需变分 $*\beta$、$\langle\cdot,\cdot\rangle_g$ 和 $\operatorname{vol}_g$。

内积变分。逆度量变化\eqnrefs{eq:math-dg-inverse-metric-variation}意味着 $k$-形式点态内积满足
\begin{equation*}
 \partial_t\langle\alpha,\beta\rangle_g
 =-\langle \mathcal V^{[k]}\alpha,\beta\rangle_g,
\end{equation*}
其中 $\mathcal V^{[k]}$ 是度量变分算子 $\mathcal V=g^{-1}\dot g$ 在 $k$-形式上的自然提升。

体积变分。由\eqnrefs{eq:math-dg-volume-variation}，
\begin{equation*}
 \partial_t\operatorname{vol}_g
 =\frac{1}{2}\operatorname{tr}_g(\dot g)\,\operatorname{vol}_g.
\end{equation*}

合并。对 Hodge 星的隐式定义等式求导，
\begin{align*}
 \alpha\wedge\partial_t(*\beta)
 &=\left[\partial_t\langle\alpha,\beta\rangle_g\right]\operatorname{vol}_g
  +\langle\alpha,\beta\rangle_g\,\partial_t\operatorname{vol}_g\\
 &=\left[
  -\langle \mathcal V^{[k]}\alpha,\beta\rangle_g
  +\frac{1}{2}\operatorname{tr}_g(\dot g)
  \langle\alpha,\beta\rangle_g
 \right]\operatorname{vol}_g.
\end{align*}

自伴性。因为 $\mathcal V$ 自伴（$g^{-1}\dot g$ 关于 $g$ 对称），诱导的 $\mathcal V^{[k]}$ 在形式内积下也自伴：
\begin{equation*}
 \langle \mathcal V^{[k]}\alpha,\beta\rangle_g
 =\langle\alpha,\mathcal V^{[k]}\beta\rangle_g.
\end{equation*}

反读 Hodge 星。将自伴性代入上面合并得到的等式，
\begin{equation*}
 \alpha\wedge\partial_t(*\beta)
 =\left\langle
  \frac{1}{2}\operatorname{tr}_g(\dot g)\,\alpha
  -\mathcal V^{[k]}\alpha,\,
  \beta
 \right\rangle_g\operatorname{vol}_g
 =\alpha\wedge *\left[
  \frac{1}{2}\operatorname{tr}_g(\dot g)\,\beta
  -\mathcal V^{[k]}\beta
 \right],
\end{equation*}
由 $\alpha$ 任意性即得到\eqnrefs{eq:math-dg-hodge-star-variation}。

共形核对：$\dot g=2ug$ 时 $\mathcal V=2uI$，
其在 $k$-形式上的提升为 $\mathcal V^{[k]}=2kuI$，故
\[
\dot *\alpha=(n-2k)u*\alpha.
\]
这与 $g\mapsto e^{2f}g$ 时 $*\mapsto e^{(n-2k)f}*$ 相容；
中间次数 $k=n/2$ 的 Hodge 星共形不变。
\end{proof}

\begin{proposition}[Ricci 孤子产生自相似流]
若 $\operatorname{Ric}_g+\frac12\mathcal L_Xg=\lambda g$，令 $c(t)=1-2\lambda t>0$，$\phi_t$ 由 $X/c(t)$ 生成。则在流存在的时间区间，$g(t)=c(t)\phi_t^*g$ 精确满足 Ricci 流方程。
\end{proposition}

\begin{proof}
设 $(g,X,\lambda)$ 满足\eqnrefs{eq:math-dg-ricci-soliton}，令
\[
c(t)=1-2\lambda t.
\]
在 $c(t)>0$ 的区间上，令 $\phi_t$ 由时间依赖向量场
\[
Y_t=\frac{1}{c(t)}X
\]
生成，并定义
\begin{equation}
g(t)=c(t)\phi_t^*g.
\label{eq:math-dg-soliton-selfsimilar-flow}
\end{equation}
则
\begin{align*}
\partial_tg(t)
&=c'(t)\phi_t^*g
+c(t)\phi_t^*(\mathcal L_{Y_t}g)\\
&=\phi_t^*(-2\lambda g+\mathcal L_Xg)\\
&=-2\phi_t^*\operatorname{Ric}_g.
\end{align*}
常数缩放不改变 $(0,2)$ 型 Ricci 张量，而 Ricci 张量 对微分同胚自然，因此
\[
\operatorname{Ric}_{g(t)}=\phi_t^*\operatorname{Ric}_g.
\]
所以
\[
\partial_tg(t)=-2\operatorname{Ric}_{g(t)}.
\]
这证明 孤子 不是"近似自相似"，而是 Ricci 流 在缩放与微分同胚作用下的精确固定点。
\end{proof}

